\chapter{Programs — Steps}
%%% generated by the blueprint pipeline from TCSlib.Programs.WordRAM.Steps
\section{Overview}

Declarations of the Lean module \texttt{TCSlib.Programs.WordRAM.Steps}.

\section{Declarations}

\begin{definition}[Signed word range]
\lean{WordRAM.InRange}
\label{WordRAM.InRange}
\leanok
An integer $v$ is in range for word width $W$ exactly when $-2^W\le 2v<2^W$.
\end{definition}

\begin{definition}[Integer encoding as a word]
\lean{WordRAM.wd}
\label{WordRAM.wd}
\leanok
For word width $W$, the word encoding of an integer $v$ is its $W$-bit representation.
\end{definition}

\begin{theorem}[Signed interpretation of an encoded integer]
\lean{WordRAM.toInt_wd}
\label{WordRAM.toInt_wd}
\leanok
\uses{WordRAM.InRange, WordRAM.wd}
\difficulty{4}
Let $W$ be a word width and let $v\in\mathbb{Z}$ satisfy $-2^W\le 2v<2^W$. The signed
integer represented by the $W$-bit encoding of $v$ is $v$.
\end{theorem}

\begin{theorem}[Encoding a bit vector’s signed value]
\lean{WordRAM.wd_toInt}
\label{WordRAM.wd_toInt}
\leanok
\uses{WordRAM.wd}
\difficulty{1}
For any word width $W$ and any $W$-bit vector $w$, encoding the signed integer
represented by $w$ gives $w$.
\end{theorem}

\begin{theorem}[Addition of encoded integers]
\lean{WordRAM.wd_add}
\label{WordRAM.wd_add}
\leanok
\uses{WordRAM.wd}
\difficulty{1}
For any word width $W$ and integers $a,b$, adding their $W$-bit encodings gives the
$W$-bit encoding of $a+b$.
\end{theorem}

\begin{theorem}[Multiplication of encoded integers]
\lean{WordRAM.wd_mul}
\label{WordRAM.wd_mul}
\leanok
\uses{WordRAM.wd}
\difficulty{1}
For any word width $W$ and integers $a,b$, multiplying their $W$-bit encodings gives the
$W$-bit encoding of $ab$.
\end{theorem}

\begin{theorem}[Subtraction of encoded integers]
\lean{WordRAM.wd_sub}
\label{WordRAM.wd_sub}
\leanok
\uses{WordRAM.wd}
\difficulty{3}
For any word width $W$ and integers $a,b$, subtracting the $W$-bit encoding of $b$ from
the $W$-bit encoding of $a$ gives the $W$-bit encoding of $a-b$.
\end{theorem}

\begin{theorem}[Encoding one]
\lean{WordRAM.wd_one}
\label{WordRAM.wd_one}
\leanok
\uses{WordRAM.wd}
\difficulty{1}
For any word width $W$, the $W$-bit representation of the bit-vector value $1$ is the
encoding of the integer $1$.
\end{theorem}

\begin{theorem}[Symmetric range from an absolute bound]
\lean{WordRAM.inRange_of_abs_le}
\label{WordRAM.inRange_of_abs_le}
\leanok
\uses{WordRAM.InRange}
\difficulty{4}
Let $v,B$ be integers and $W$ a natural number. If $2B<2^W$ and $|v|\le B$, then
$-2^W\le 2v<2^W$.
\end{theorem}

\begin{theorem}[Loaded word at a natural address]
\lean{WordRAM.loadWords_mem_nat}
\label{WordRAM.loadWords_mem_nat}
\leanok
\uses{WordRAM.loadWords, WordRAM.wd}
\difficulty{3}
For every natural word width $W$, list of integers $ws$, and natural address $a$, the
memory obtained by loading $ws$ at address $a$ contains at that address the $W$-bit
representation of the $a$th list entry, using zero when the list has no such entry.
\end{theorem}

\begin{theorem}[Loaded word at a negative address]
\lean{WordRAM.loadWords_mem_neg}
\label{WordRAM.loadWords_mem_neg}
\leanok
\uses{WordRAM.loadWords, WordRAM.wd}
\difficulty{1}
For every natural word width $W$, integer list $ws$, and negative integer address $a$,
the loaded memory contains the zero word at $a$.
\end{theorem}

\begin{definition}[Machine configuration]
\lean{WordRAM.Cfg}
\label{WordRAM.Cfg}
\leanok
For word width $W$, a machine configuration consists of a natural-number program counter
and a memory assigning a $W$-bit word to each integer address.
\end{definition}

\begin{definition}[One machine step]
\lean{WordRAM.step}
\label{WordRAM.step}
\leanok
\uses{WordRAM.Cfg, WordRAM.Instr}
For a program $P$ and a configuration with program counter $p$ and memory $m$, one step
returns either a new configuration or a Boolean terminal result. If the instruction at
$p$ is missing, it is treated as rejection. An instruction that writes a value to
address $i$ stores that $W$-bit value at $i$ and advances the counter to $p+1$; the
arithmetic instructions write the corresponding word addition, subtraction, or
multiplication, and load and store use the integer value represented by a word as an
address. A negative-value branch jumps to its specified position exactly when the tested
word represents a negative integer, and otherwise advances to $p+1$; accept and reject
return true and false, respectively.
\end{definition}

\begin{definition}[Iterated machine steps]
\lean{WordRAM.Steps}
\label{WordRAM.Steps}
\leanok
\uses{WordRAM.Cfg, WordRAM.Instr, WordRAM.step}
For a program $P$ and configurations $c,c'$, zero steps relate them exactly when $c=c'$.
A run of $n+1$ steps from $c$ to $c'$ exists exactly when one step from $c$ returns a
configuration $c_1$ and there is a run of $n$ steps from $c_1$ to $c'$.
\end{definition}

\begin{theorem}[Reflexivity of zero-step execution]
\lean{WordRAM.Steps.refl}
\label{WordRAM.Steps.refl}
\leanok
\uses{WordRAM.Cfg, WordRAM.Steps}
\difficulty{1}
For any word size $W$, program $P$, and configuration $c$ with that word size, zero
steps of $P$ take $c$ to itself.
\end{theorem}

\begin{theorem}[A Single Step as a One-Step Run]
\lean{WordRAM.Steps.one}
\label{WordRAM.Steps.one}
\leanok
\uses{WordRAM.Cfg, WordRAM.Steps, WordRAM.step}
\difficulty{1}
For any program and configurations $c,c'$, if one step of the program starting at $c$
produces $c'$ rather than halting, then one step takes $c$ to $c'$.
\end{theorem}

\begin{theorem}[Concatenation of machine runs]
\lean{WordRAM.Steps.trans}
\label{WordRAM.Steps.trans}
\leanok
\uses{WordRAM.Cfg, WordRAM.Steps}
\difficulty{4}
For a WordRAM program, if a configuration $c$ reaches $c_1$ in $n$ steps and $c_1$
reaches $c_2$ in $k$ steps, then $c$ reaches $c_2$ in $n+k$ steps.
\end{theorem}

\begin{theorem}[Changing the Final Program Counter]
\lean{WordRAM.Steps.cast_pos}
\label{WordRAM.Steps.cast_pos}
\leanok
\uses{WordRAM.Cfg, WordRAM.Steps}
\difficulty{1}
For any program, natural number $n$, program counters $p,p'$, initial configuration $c$,
and memory function $m$, if $n$ steps take $c$ to the configuration with program counter
$p$ and memory $m$, and $p=p'$, then $n$ steps take $c$ to the configuration with
program counter $p'$ and memory $m$.
\end{theorem}

\begin{theorem}[Changing the Step Count]
\lean{WordRAM.Steps.cast_count}
\label{WordRAM.Steps.cast_count}
\leanok
\uses{WordRAM.Cfg, WordRAM.Steps}
\difficulty{1}
For any program, natural numbers $n,n'$, and configurations $c,c'$, if $n$ steps of the
program take $c$ to $c'$ and $n=n'$, then $n'$ steps take $c$ to $c'$.
\end{theorem}

\begin{theorem}[Zero-step execution]
\lean{WordRAM.exec_zero}
\label{WordRAM.exec_zero}
\leanok
\uses{WordRAM.Cfg, WordRAM.exec}
\difficulty{1}
For every natural word width $W$, every program $P$ consisting of WordRAM instructions,
and every configuration $c$ of width $W$, executing $P$ for zero steps from the program
counter and memory of $c$ returns no result.
\end{theorem}

\begin{theorem}[One-step execution equation]
\lean{WordRAM.exec_succ}
\label{WordRAM.exec_succ}
\leanok
\uses{WordRAM.Cfg, WordRAM.exec, WordRAM.step}
\difficulty{4}
Let $W$ be a word width, $P$ a finite program, $t$ a nonnegative integer, and $c$ a
configuration consisting of a program counter and an integer-indexed memory of $W$-bit
words. Executing $P$ from $c$ for $t+1$ steps returns the halted verdict with $c$'s
memory if the next step halts; if that step continues at a configuration $c_1$, it
returns the result of executing $P$ from $c_1$ for $t$ steps.
\end{theorem}

\begin{theorem}[Execution after a halting step]
\lean{WordRAM.exec_succ_of_verdict}
\label{WordRAM.exec_succ_of_verdict}
\leanok
\uses{WordRAM.Cfg, WordRAM.exec, WordRAM.exec_succ, WordRAM.step}
\difficulty{1}
Let $W$ be a word width, $P$ a finite program, and $c$ a configuration consisting of a
program counter and an integer-indexed memory of $W$-bit words. If one step from $c$
halts with Boolean verdict $v$, then for every nonnegative integer $t$, executing $P$
from $c$ for $t+1$ steps returns $v$ and leaves the memory unchanged.
\end{theorem}

\begin{theorem}[Execution after a continuing step]
\lean{WordRAM.exec_succ_of_step}
\label{WordRAM.exec_succ_of_step}
\leanok
\uses{WordRAM.Cfg, WordRAM.exec, WordRAM.exec_succ, WordRAM.step}
\difficulty{1}
Let $W$ be a word width, $P$ a finite program, and $c,c'$ configurations consisting of a
program counter and an integer-indexed memory of $W$-bit words. If one step from $c$
continues at $c'$, then for every nonnegative integer $t$, executing $P$ from $c$ for
$t+1$ steps gives the same result as executing it from $c'$ for $t$ steps.
\end{theorem}

\begin{theorem}[Execution from a sequence of steps]
\lean{WordRAM.exec_of_steps}
\label{WordRAM.exec_of_steps}
\leanok
\uses{WordRAM.Cfg, WordRAM.Steps, WordRAM.exec, WordRAM.exec_succ_of_step, WordRAM.exec_succ_of_verdict, WordRAM.step}
\difficulty{4}
Let $W$ be a word width, $P$ a finite program, $n$ a nonnegative integer, and $c,c'$
configurations consisting of a program counter and an integer-indexed memory of $W$-bit
words. Suppose $c$ reaches $c'$ in exactly $n$ steps, each of which continues to another
configuration, and the next step from $c'$ halts with Boolean verdict $v$. Then
executing $P$ from the program counter and memory of $c$ for $n+1$ steps returns verdict
$v$ and the memory of $c'$.
\end{theorem}

\begin{theorem}[Monotonicity of bounded execution]
\lean{WordRAM.exec_mono}
\label{WordRAM.exec_mono}
\leanok
\uses{WordRAM.Cfg, WordRAM.exec, WordRAM.exec_succ_of_step, WordRAM.exec_succ_of_verdict, WordRAM.exec_zero, WordRAM.step}
\difficulty{5}
Let $W$ be a word width, $P$ a finite program, $c$ a configuration consisting of a
program counter and a memory of $W$-bit words indexed by integers, and $r$ a pair
consisting of a Boolean verdict and a memory. If executing $P$ from $c$ for $t$ steps
returns $r$, then for every $t' \ge t$, executing $P$ from the same configuration for
$t'$ steps also returns $r$.
\end{theorem}

\begin{definition}[Code at a Position]
\lean{WordRAM.CodeAt}
\label{WordRAM.CodeAt}
\leanok
\uses{WordRAM.Instr}
For an instruction list $P$, a position $p$, and a list of instructions $l$, say that
$l$ is code at position $p$ in $P$ when $l$ is a prefix of the list obtained by removing
the first $p$ instructions of $P$.
\end{definition}

\begin{theorem}[Instruction sequence at its initial position]
\lean{WordRAM.codeAt_self}
\label{WordRAM.codeAt_self}
\leanok
\uses{WordRAM.CodeAt, WordRAM.Instr}
\difficulty{1}
Every finite instruction sequence is the initial segment of itself when read from
position zero.
\end{theorem}

\begin{theorem}[Left part of a code segment]
\lean{WordRAM.CodeAt.left}
\label{WordRAM.CodeAt.left}
\leanok
\uses{WordRAM.CodeAt, WordRAM.Instr}
\difficulty{1}
Let $P$ be a list of instructions, and suppose the concatenation $l_1\mathbin{++}l_2$ is
a prefix of the part of $P$ beginning at position $\mathrm{pos}$. Then $l_1$ is also a
prefix of that part of $P$.
\end{theorem}

\begin{theorem}[Right part of a code segment]
\lean{WordRAM.CodeAt.right}
\label{WordRAM.CodeAt.right}
\leanok
\uses{WordRAM.CodeAt, WordRAM.Instr}
\difficulty{3}
Let $P$ be a list of instructions, and suppose $l_1\mathbin{++}l_2$ is a prefix of the
part of $P$ beginning at position $\mathrm{pos}$. Then $l_2$ is a prefix of the part
beginning at position $\mathrm{pos}+|l_1|$, where $|l_1|$ is the length of $l_1$.
\end{theorem}

\begin{theorem}[First instruction at a code position]
\lean{WordRAM.CodeAt.head}
\label{WordRAM.CodeAt.head}
\leanok
\uses{WordRAM.CodeAt, WordRAM.Instr}
\difficulty{3}
Let $P$ be a list of instructions, and suppose the list beginning at position
$\mathrm{pos}$ has $i$ as its first instruction. Then the instruction at position
$\mathrm{pos}$ in $P$ is $i$; when that position is beyond the end of $P$, the
instruction there is taken to be reject.
\end{theorem}

\begin{theorem}[Tail of a code segment]
\lean{WordRAM.CodeAt.tail}
\label{WordRAM.CodeAt.tail}
\leanok
\uses{WordRAM.CodeAt, WordRAM.CodeAt.right, WordRAM.Instr}
\difficulty{1}
Let $P$ be a list of instructions, and suppose the list beginning at position
$\mathrm{pos}$ has $i$ followed by the list $l$. Then $l$ is a prefix of the list
beginning at position $\mathrm{pos}+1$.
\end{theorem}

\begin{theorem}[Changing the code position notation]
\lean{WordRAM.CodeAt.cast_pos}
\label{WordRAM.CodeAt.cast_pos}
\leanok
\uses{WordRAM.CodeAt, WordRAM.Instr}
\difficulty{1}
Let $P$ be a list of instructions, let $l$ be a list of instructions that is a prefix of
the part of $P$ beginning at position $\mathrm{pos}$, and suppose
$\mathrm{pos}=\mathrm{pos}'$. Then $l$ is a prefix of the part of $P$ beginning at
position $\mathrm{pos}'$.
\end{theorem}

\begin{definition}[Straight-Line Instructions]
\lean{WordRAM.Straight}
\label{WordRAM.Straight}
\leanok
\uses{WordRAM.Instr}
An instruction is straight-line if it is an assignment of the constant $1$, an addition,
subtraction, or multiplication of register values, a load from memory, or a store to
memory.
\end{definition}

\begin{definition}[Effect of an Instruction]
\lean{WordRAM.effect}
\label{WordRAM.effect}
\leanok
\uses{WordRAM.Instr, WordRAM.wd}
For a word width $W$, an instruction acts on a memory state $m$ by updating the
destination register for a constant assignment, arithmetic operation, or load, or the
addressed memory cell for a store. A constant assignment writes the $W$-bit word
representing $1$; arithmetic uses word arithmetic, and loads and stores use the integer
represented by the relevant word as the memory address. All other instructions leave the
memory state unchanged.
\end{definition}

\begin{definition}[Effect of an Instruction List]
\lean{WordRAM.effects}
\label{WordRAM.effects}
\leanok
\uses{WordRAM.Instr, WordRAM.effect}
For a word width $W$, the effect of a list of instructions on a memory state is obtained
by applying the effect of each instruction in the list successively, from first to last.
\end{definition}

\begin{theorem}[Effect of the empty instruction list]
\lean{WordRAM.effects_nil}
\label{WordRAM.effects_nil}
\leanok
\uses{WordRAM.effects}
\difficulty{1}
For any word width $W$ and any memory $m$ assigning a $W$-bit word to each integer
address, applying the empty instruction list leaves $m$ unchanged.
\end{theorem}

\begin{theorem}[Effect of a nonempty instruction list]
\lean{WordRAM.effects_cons}
\label{WordRAM.effects_cons}
\leanok
\uses{WordRAM.Instr, WordRAM.effect, WordRAM.effects}
\difficulty{1}
Fix a word width $W$, an instruction $i$, a finite instruction list $l$, and an initial
memory $m$ assigning a $W$-bit word to each integer address. Applying the list beginning
with $i$ to $m$ gives the same memory as first applying $i$ and then applying $l$.
\end{theorem}

\begin{theorem}[Composition of instruction effects]
\lean{WordRAM.effects_append}
\label{WordRAM.effects_append}
\leanok
\uses{WordRAM.Instr, WordRAM.effects}
\difficulty{3}
Fix a word width $W$, a program instruction type, two finite instruction lists
$l_1,l_2$, and a memory $m$ assigning a $W$-bit word to each integer address. Applying
the instructions in the concatenated list $l_1 ++ l_2$ to $m$ gives the same memory as
first applying $l_1$ and then applying $l_2$.
\end{theorem}

\begin{theorem}[One step for a straight instruction]
\lean{WordRAM.step_straight}
\label{WordRAM.step_straight}
\leanok
\uses{WordRAM.Cfg, WordRAM.Instr, WordRAM.Straight, WordRAM.effect, WordRAM.step, WordRAM.wd_one}
\difficulty{4}
Let $W$ be a word width, $P$ a program, $c$ a machine configuration, and $i$ an
instruction. If the instruction at $c$'s program counter is $i$ and $i$ is a straight
instruction (one, addition, subtraction, multiplication, load, or store), then one step
of $P$ from $c$ advances the program counter by one and changes the memory to the effect
of $i$ on $c$'s memory.
\end{theorem}

\begin{theorem}[One step for a straight instruction]
\lean{WordRAM.steps_instr}
\label{WordRAM.steps_instr}
\leanok
\uses{WordRAM.CodeAt, WordRAM.CodeAt.head, WordRAM.Instr, WordRAM.Steps, WordRAM.Steps.one, WordRAM.Straight, WordRAM.effect, WordRAM.step_straight}
\difficulty{1}
Let $W$ be a word width, $P$ a program, $pos$ a program counter, $m$ a memory, and $i$
an instruction. If the program at $pos$ begins with $i$ and $i$ is a straight
instruction (one, addition, subtraction, multiplication, load, or store), then one step
from $(pos,m)$ reaches program counter $pos+1$ with memory equal to the effect of $i$ on
$m$.
\end{theorem}

\begin{theorem}[Execution of a straight code segment]
\lean{WordRAM.steps_straight}
\label{WordRAM.steps_straight}
\leanok
\uses{WordRAM.CodeAt, WordRAM.CodeAt.head, WordRAM.CodeAt.tail, WordRAM.Instr, WordRAM.Steps, WordRAM.Steps.cast_pos, WordRAM.Straight, WordRAM.effect, WordRAM.effects, WordRAM.step_straight}
\difficulty{4}
Let $W$ be a word width, $P$ a program, $l$ a list of instructions, $pc$ a program
counter, and $m$ a memory. If $l$ is a prefix of the program beginning at $pc$ and every
instruction in $l$ is straight (one, addition, subtraction, multiplication, load, or
store), then executing $|l|$ steps from $(pc,m)$ reaches program counter $pc+|l|$ with
memory obtained by applying the effects of the instructions in $l$ successively to $m$.
\end{theorem}

\begin{theorem}[Effect of Writing One]
\lean{WordRAM.effect_one}
\label{WordRAM.effect_one}
\leanok
\uses{WordRAM.effect, WordRAM.wd}
\difficulty{1}
Let $W\in\mathbb{N}$, let $m:\mathbb{Z}\to\{0,1\}^W$ be a memory, and let
$i\in\mathbb{Z}$. Executing the instruction that writes one at address $i$ changes the
memory to the function obtained by updating address $i$ with the $W$-bit encoding of
$1$.
\end{theorem}

\begin{theorem}[Effect of Addition]
\lean{WordRAM.effect_add}
\label{WordRAM.effect_add}
\leanok
\uses{WordRAM.effect, WordRAM.wd, WordRAM.wd_add}
\difficulty{3}
Let $W\in\mathbb{N}$, let $m:\mathbb{Z}\to\{0,1\}^W$ be a memory, and let
$i,j,k,a,b\in\mathbb{Z}$. If the words stored at addresses $j$ and $k$ are the $W$-bit
encodings of $a$ and $b$, respectively, then executing the addition instruction with
destination $i$ changes the memory to the function obtained by updating address $i$ with
the $W$-bit encoding of $a+b$.
\end{theorem}

\begin{theorem}[Effect of Subtraction]
\lean{WordRAM.effect_sub}
\label{WordRAM.effect_sub}
\leanok
\uses{WordRAM.effect, WordRAM.wd, WordRAM.wd_sub}
\difficulty{3}
Let $W\in\mathbb{N}$, let $m:\mathbb{Z}\to\{0,1\}^W$ be a memory, and let
$i,j,k,a,b\in\mathbb{Z}$. If the words stored at addresses $j$ and $k$ are the $W$-bit
encodings of $a$ and $b$, respectively, then executing the subtraction instruction with
destination $i$ changes the memory to the function obtained by updating address $i$ with
the $W$-bit encoding of $a-b$.
\end{theorem}

\begin{theorem}[Effect of Multiplication]
\lean{WordRAM.effect_mul}
\label{WordRAM.effect_mul}
\leanok
\uses{WordRAM.effect, WordRAM.wd, WordRAM.wd_mul}
\difficulty{3}
Let $W\in\mathbb{N}$, let $m:\mathbb{Z}\to\{0,1\}^W$ be a memory, and let
$i,j,k,a,b\in\mathbb{Z}$. If the words stored at addresses $j$ and $k$ are the $W$-bit
encodings of $a$ and $b$, respectively, then executing the multiplication instruction
with destination $i$ changes the memory to the function obtained by updating address $i$
with the $W$-bit encoding of $ab$.
\end{theorem}

\begin{theorem}[Effect of Loading]
\lean{WordRAM.effect_load}
\label{WordRAM.effect_load}
\leanok
\uses{WordRAM.InRange, WordRAM.effect, WordRAM.toInt_wd, WordRAM.wd}
\difficulty{3}
Let $W\in\mathbb{N}$, let $m:\mathbb{Z}\to\{0,1\}^W$ be a memory, and let
$i,j,a\in\mathbb{Z}$. Suppose the word at address $j$ is the $W$-bit encoding of $a$,
and $a$ satisfies $-2^W\le 2a<2^W$. Then executing the load instruction with destination
address $i$ and pointer address $j$ changes the memory to the function obtained by
updating address $i$ with the word stored at address $a$.
\end{theorem}

\begin{theorem}[Effect of Storing]
\lean{WordRAM.effect_store}
\label{WordRAM.effect_store}
\leanok
\uses{WordRAM.InRange, WordRAM.effect, WordRAM.toInt_wd, WordRAM.wd}
\difficulty{3}
Let $W\in\mathbb{N}$, let $m:\mathbb{Z}\to\{0,1\}^W$ be a memory, and let
$i,j,a\in\mathbb{Z}$. Suppose the word at address $i$ is the $W$-bit encoding of $a$,
and $a$ satisfies $-2^W\le 2a<2^W$. Then executing the store instruction with pointer
address $i$ and source address $j$ changes the memory to the function obtained by
updating address $a$ with the word stored at address $j$.
\end{theorem}

\begin{theorem}[One-step addition]
\lean{WordRAM.steps_add}
\label{WordRAM.steps_add}
\leanok
\uses{WordRAM.CodeAt, WordRAM.Instr, WordRAM.Steps, WordRAM.effect_add, WordRAM.steps_instr, WordRAM.wd}
\difficulty{1}
Let $W$ be a word width, $P$ a program, $m$ a memory, and $i,j,k,a,b$ integers. If the
program beginning at position $pos$ starts with the instruction that adds the contents
of addresses $j$ and $k$ into address $i$, and the words at $j$ and $k$ represent $a$
and $b$, respectively, then one step from $(pos,m)$ reaches program counter $pos+1$ and
memory equal to $m$ updated at address $i$ with the word representing $a+b$.
\end{theorem}

\begin{theorem}[One subtraction step]
\lean{WordRAM.steps_sub}
\label{WordRAM.steps_sub}
\leanok
\uses{WordRAM.CodeAt, WordRAM.Instr, WordRAM.Steps, WordRAM.effect_sub, WordRAM.steps_instr, WordRAM.wd}
\difficulty{1}
Let $W$ be a word width, let $P$ be a program, and let $m:\mathbb{Z}\to\{0,1\}^W$ be a
memory. Suppose the program beginning at position $p$ has a prefix whose first
instruction is subtraction of the contents of registers $j$ and $k$ into register $i$,
and suppose $m(j)$ and $m(k)$ are the $W$-bit encodings of integers $a$ and $b$. Then
one step from position $p$ yields position $p+1$ and memory obtained by replacing
register $i$ with the $W$-bit encoding of $a-b$.
\end{theorem}

\begin{theorem}[One-step load]
\lean{WordRAM.steps_load}
\label{WordRAM.steps_load}
\leanok
\uses{WordRAM.CodeAt, WordRAM.InRange, WordRAM.Instr, WordRAM.Steps, WordRAM.effect_load, WordRAM.steps_instr, WordRAM.wd}
\difficulty{1}
Let $W$ be a word width, $P$ a program, $m$ a memory, and $i,j,a$ integers. Suppose the
program at position $pos$ begins with the instruction that loads into address $i$ the
word at the address stored at $j$, the word at $j$ represents $a$, and $a$ satisfies
$-2^W \le 2a < 2^W$. Then one step from $(pos,m)$ reaches program counter $pos+1$ and
memory obtained by updating $m$ at address $i$ with the word $m(a)$.
\end{theorem}

\begin{theorem}[One-step store]
\lean{WordRAM.steps_store}
\label{WordRAM.steps_store}
\leanok
\uses{WordRAM.CodeAt, WordRAM.InRange, WordRAM.Instr, WordRAM.Steps, WordRAM.effect_store, WordRAM.steps_instr, WordRAM.wd}
\difficulty{1}
Let $W$ be a word width, $P$ a program, $m$ a memory, and $i,j,a$ integers. Suppose the
program at position $pos$ begins with the instruction that stores the word at address
$j$ at the address held in $i$, the word at $i$ represents $a$, and $a$ satisfies $-2^W
\le 2a < 2^W$. Then one step from $(pos,m)$ reaches program counter $pos+1$ and memory
obtained by updating $m$ at address $a$ with the word $m(j)$.
\end{theorem}

\begin{theorem}[One-step conditional branch]
\lean{WordRAM.steps_bltz}
\label{WordRAM.steps_bltz}
\leanok
\uses{WordRAM.CodeAt, WordRAM.CodeAt.head, WordRAM.InRange, WordRAM.Instr, WordRAM.Steps, WordRAM.Steps.one, WordRAM.step, WordRAM.toInt_wd, WordRAM.wd}
\difficulty{1}
Let $W$ be a word width, $P$ a program, and $c,l,pos$ integers or natural numbers as
appropriate, with $v$ an integer and $m$ a memory. Suppose the program at position $pos$
begins with the instruction that branches to $l$ when address $c$ contains a negative
value, the word at $c$ represents $v$, and $v$ satisfies $-2^W \le 2v < 2^W$. Then one
step from $(pos,m)$ leaves the memory unchanged and moves the program counter to $l$ if
$v<0$, and to $pos+1$ otherwise.
\end{theorem}

\begin{theorem}[Branching and Step Sequences]
\lean{WordRAM.Steps.branch}
\label{WordRAM.Steps.branch}
\leanok
\uses{WordRAM.Cfg, WordRAM.Steps}
\difficulty{4}
Let $n,l,pos,pos'$ be natural numbers, let $c$ be an initial configuration, let $m$ be a
memory function, and let $v$ be an integer. Suppose that $n$ steps of the program take
$c$ to the configuration with memory $m$ and program counter $l$ if $v<0$, and with
memory $m$ and program counter $pos$ otherwise. If a proposition $p$ is equivalent to
$v\ge 0$ and $pos=pos'$, then if $p$ holds, $n$ steps take $c$ to the configuration with
memory $m$ and program counter $pos'$; if $p$ fails, they take $c$ to the configuration
with memory $m$ and program counter $l$.
\end{theorem}

\begin{theorem}[Accept instruction halts with acceptance]
\lean{WordRAM.step_accept}
\label{WordRAM.step_accept}
\leanok
\uses{WordRAM.Cfg, WordRAM.step}
\difficulty{1}
Let $P$ be a Word RAM program and $c$ a configuration. If the instruction at the program
counter of $c$, using rejection as the default instruction when the counter is out of
range, is accept, then one step of $P$ from $c$ returns the acceptance result.
\end{theorem}

\begin{theorem}[Reject instruction halts with rejection]
\lean{WordRAM.step_reject}
\label{WordRAM.step_reject}
\leanok
\uses{WordRAM.Cfg, WordRAM.step}
\difficulty{1}
Let $P$ be a Word RAM program and $c$ a configuration. If the instruction at the program
counter of $c$, using rejection as the default instruction when the counter is out of
range, is reject, then one step of $P$ from $c$ returns the rejection result.
\end{theorem}

\begin{definition}[Agreement Outside a Set]
\lean{WordRAM.AgreeOutside}
\label{WordRAM.AgreeOutside}
\leanok
Two memory states $m$ and $m'$ agree outside a set $S$ of integer addresses if they have
equal values at every address outside $S$.
\end{definition}

\begin{theorem}[Reflexive agreement outside a set]
\lean{WordRAM.AgreeOutside.refl}
\label{WordRAM.AgreeOutside.refl}
\leanok
\uses{WordRAM.AgreeOutside}
\difficulty{1}
For every set $s$ of integer addresses and every memory $m$ mapping integer addresses to
$W$-bit words, $m$ agrees with itself outside $s$.
\end{theorem}

\begin{theorem}[Transitivity of outside agreement]
\lean{WordRAM.AgreeOutside.trans}
\label{WordRAM.AgreeOutside.trans}
\leanok
\uses{WordRAM.AgreeOutside}
\difficulty{1}
Let $s$ be a set of integer addresses, and let $m_1,m_2,m_3$ be memories mapping integer
addresses to $W$-bit words. If $m_2$ agrees with $m_1$ outside $s$ and $m_3$ agrees with
$m_2$ outside $s$, then $m_3$ agrees with $m_1$ outside $s$.
\end{theorem}

\begin{theorem}[Agreement outside a larger set]
\lean{WordRAM.AgreeOutside.mono}
\label{WordRAM.AgreeOutside.mono}
\leanok
\uses{WordRAM.AgreeOutside}
\difficulty{1}
Let $s\subseteq t$ be sets of integer addresses, and let $m,m'$ be memories mapping
integer addresses to $W$-bit words. If $m'$ agrees with $m$ outside $s$, then it agrees
with $m$ outside $t$.
\end{theorem}

\begin{theorem}[Transitivity over a union]
\lean{WordRAM.AgreeOutside.trans_union}
\label{WordRAM.AgreeOutside.trans_union}
\leanok
\uses{WordRAM.AgreeOutside, WordRAM.AgreeOutside.mono, WordRAM.AgreeOutside.trans}
\difficulty{3}
Let $s,t$ be sets of integer addresses, and let $m_1,m_2,m_3$ be memories mapping
integer addresses to $W$-bit words. If $m_2$ agrees with $m_1$ outside $s$ and $m_3$
agrees with $m_2$ outside $t$, then $m_3$ agrees with $m_1$ outside $s\cup t$.
\end{theorem}

\begin{theorem}[Agreement at an outside cell]
\lean{WordRAM.AgreeOutside.cell}
\label{WordRAM.AgreeOutside.cell}
\leanok
\uses{WordRAM.AgreeOutside}
\difficulty{1}
Let $s$ be a set of integer addresses and let $m,m'$ be memories mapping integer
addresses to $W$-bit words. If $m'$ agrees with $m$ at every address outside $s$, then
$m'(a)=m(a)$ for every address $a$ outside $s$.
\end{theorem}

\begin{theorem}[Reading an unchanged cell]
\lean{WordRAM.AgreeOutside.read}
\label{WordRAM.AgreeOutside.read}
\leanok
\uses{WordRAM.AgreeOutside}
\difficulty{1}
Let $s$ be a set of integer addresses and let $m,m'$ be memories mapping integer
addresses to $W$-bit words. If $m'$ agrees with $m$ outside $s$, if $a\notin s$, and if
$m(a)=v$ for a $W$-bit word $v$, then $m'(a)=v$.
\end{theorem}

\begin{theorem}[Updating a memory outside a set]
\lean{WordRAM.AgreeOutside.update}
\label{WordRAM.AgreeOutside.update}
\leanok
\uses{WordRAM.AgreeOutside}
\difficulty{3}
Let $s\subseteq\mathbb Z$, let $m,m':\mathbb Z\to\mathrm{BitVec}(W)$ agree at every
address outside $s$, and let $a\in s$. For any word $v\in\mathrm{BitVec}(W)$, replacing
the value of $m'$ at $a$ by $v$ gives a memory that still agrees with $m$ outside $s$.
\end{theorem}

\begin{definition}[Writes Within a Set]
\lean{WordRAM.WritesIn}
\label{WordRAM.WritesIn}
\leanok
\uses{WordRAM.Instr}
Given a set $S$ of integer addresses, an instruction writes in $S$ if it is a constant
assignment, arithmetic operation, or load whose destination register belongs to $S$. A
store is not considered to write in $S$, while every other instruction is considered to
write in $S$.
\end{definition}

\begin{theorem}[Monotonicity of write-address restrictions]
\lean{WordRAM.WritesIn.mono}
\label{WordRAM.WritesIn.mono}
\leanok
\uses{WordRAM.Instr, WordRAM.WritesIn}
\difficulty{4}
For any WordRAM instruction, if every address it writes (when it writes to an address)
lies in a set $s\subseteq\mathbb Z$, then every address it writes also lies in any set
$t$ containing $s$.
\end{theorem}

\begin{theorem}[Memory agreement outside a write set]
\lean{WordRAM.agreeOutside_effect}
\label{WordRAM.agreeOutside_effect}
\leanok
\uses{WordRAM.AgreeOutside, WordRAM.AgreeOutside.refl, WordRAM.AgreeOutside.update, WordRAM.Instr, WordRAM.WritesIn, WordRAM.effect}
\difficulty{4}
Let $W$ be a word width, $s\subseteq\mathbb Z$ a set of memory addresses, $i$ an
instruction, and $m$ a memory assigning a $W$-bit word to each integer address. If the
instruction's write location belongs to $s$ according to the condition defining writes
in $s$, then applying its memory effect to $m$ leaves $m$ unchanged at every address
outside $s$.
\end{theorem}

\begin{theorem}[Memory agreement after a sequence of effects]
\lean{WordRAM.agreeOutside_effects}
\label{WordRAM.agreeOutside_effects}
\leanok
\uses{WordRAM.AgreeOutside, WordRAM.AgreeOutside.refl, WordRAM.AgreeOutside.trans, WordRAM.Instr, WordRAM.WritesIn, WordRAM.agreeOutside_effect, WordRAM.effects}
\difficulty{4}
Let $W$ be a word width, $s\subseteq\mathbb Z$ a set of memory addresses, $l$ a finite
list of instructions, and $m$ a memory assigning a $W$-bit word to each integer address.
If every instruction in $l$ satisfies the condition that it writes in $s$, then applying
the effects of the instructions in order leaves $m$ unchanged at every address outside
$s$.
\end{theorem}
